\documentclass{article}
\usepackage[utf8]{inputenc}
\usepackage[T2A]{fontenc}
\usepackage{helvet}
\usepackage{geometry}
\usepackage{xcolor}
\usepackage{eso-pic}
\usepackage{fancyhdr}
\usepackage{parskip}
\usepackage{microtype}
\usepackage[english, ukrainian]{babel}
\usepackage{url}
\usepackage{amsmath}
\usepackage{amssymb}
\usepackage{amsthm}
\usepackage{longtable}
\usepackage{booktabs}
\usepackage{array}
\usepackage{hyperref}
\usepackage{titlesec}
\usepackage{tikz}
\usetikzlibrary{positioning, shapes.geometric, arrows.meta}

% Define brand colors
\definecolor{primary}{HTML}{293757}
\definecolor{footergray}{gray}{0.65}

\geometry{a4paper,left=3cm,right=3cm,top=3cm,bottom=3cm}
\renewcommand{\familydefault}{\sfdefault}
\setcounter{secnumdepth}{3}

% Font shapes for T2A/Helvetica (mapping to cmss for Cyrillic support)
\DeclareFontFamily{T2A}{phv}{}
\DeclareFontShape{T2A}{phv}{m}{n}{<->ssub * cmss/m/n}{}
\DeclareFontShape{T2A}{phv}{bx}{n}{<->ssub * cmss/bx/n}{}
\DeclareFontShape{T2A}{phv}{m}{it}{<->ssub * cmss/m/it}{}
\DeclareFontShape{T2A}{phv}{bx}{it}{<->ssub * cmss/bx/it}{}

\titleformat{\section}
  {\color{primary}\Huge\bfseries}{\thesection.}{1em}{}
\titlespacing*{\section}{0pt}{30pt}{12pt}
\titleformat{\subsection}
  {\color{primary}\LARGE\bfseries}{\thesubsection.}{1em}{}
\titlespacing*{\subsection}{0pt}{20pt}{8pt}
\titleformat{\subsubsection}
  {\color{primary}\Large\bfseries}{\thesubsubsection.}{1em}{}
\titlespacing*{\subsubsection}{0pt}{10pt}{6pt}

\fancyhf{}
\fancyhead[R]{\small\color{footergray} 29 червня 2026}
\fancyfoot[L]{\small\color{footergray} Техноробочий проект ERP/1. Модуль Комунікатор}
\fancyfoot[R]{\small\color{footergray}\thepage}

\newcommand\HUGE[1]{\fontsize{38}{38}\selectfont #1}
\renewcommand{\headrulewidth}{0pt}
\renewcommand{\footrulewidth}{0.4pt}
\renewcommand{\footrule}{\color{footergray}\hrule width \textwidth height 0.4pt}

\begin{document}
\pagestyle{fancy}

\begin{titlepage}
\AddToShipoutPictureBG*{\AtPageLowerLeft{\color{primary}\rule{\paperwidth}{\paperheight}}}
\color{white}\thispagestyle{empty}\vspace*{3cm}\centering
  {\Huge \textbf{ERP/1: Комунікатор} \par} \vspace{1.5cm}
  {\HUGE \textbf{Техноробочий проект} \par} \vspace{1.5cm}
  {\LARGE \textbf{на створення та впровадження \\ локальної системи безпечного зв'язку, \\ видачі сертифікатів та обміну повідомленнями} \par} \vspace{1.5cm}
  {\large Згідно з Постановою КМУ № 205 від 21.02.2025 \par}
  {\large Формалізація вимог за стандартом ISO/IEC/IEEE 29148:2018 \par}
  \vfill
  {\large 2026 \copyright\ Системи Електронного Урядування \par}
\end{titlepage}

\newpage
\begin{center}
  {\Huge\color{primary}\bfseries Зміст\par}
\end{center}
\vspace{40pt}
\tableofcontents

\newpage
\section{Загальні відомості про засіб інформатизації}

\subsection{Найменування та підстави для розробки}
\textbf{Найменування засобу інформатизації:} Модуль «Комунікатор» (інфраструктура безпечного зв'язку CHAT v2 / MAIL v3) у складі інформаційної системи управління підприємством ERP/1.

\textbf{Підстави для виконання робіт:}
\begin{itemize}
    \item Закон України «Про захист персональних даних»;
    \item Закон України «Про захист інформації в інформаційно-комунікаційних системах»;
    \item Порядок використання засобів інформатизації (Постанова КМУ № 205 від 21.02.2025);
    \item Регламент взаємодії систем безпеки та шифрування X.422.
\end{itemize}

\subsection{Призначення та цілі створення засобу}
Модуль призначений для розгортання периферійних систем захищеного обміну повідомленнями та документами, локального управління сертифікатами відкритих ключів (CA), VPN-тунелювання, служби каталогів (LDAP) та Identity Access Management (IAS) без використання сторонніх хмарних рішень та мови Rust на периферії.

\subsection{План-графік виконання етапів робіт}
\begin{enumerate}
    \item \textbf{Етап 1: Технічне проектування} --- опис ASN.1 DER структур, схем каталогів LDAP та специфікацій WireGuard тунелів. (1.5 місяці).
    \item \textbf{Етап 2: Реалізація CA та EST/CMP} --- запуск локального засвідчувального центру для видачі та перевірки статусів X.509. (2 місяці).
    \item \textbf{Етап 3: Розробка CHAT v2 та MAIL v3 серверів} --- впровадження асинхронного брокера повідомлень на базі Erlang/OTP та поштової черги доставки документів. (2.5 місяці).
    \item \textbf{Етап 4: Інтеграція клієнтських додатків} --- реалізація клієнтських Swift/SwiftUI SDK для роботи за протоколом Buddha X.422. (1.5 місяці).
    \item \textbf{Етап 5: Державна експертиза КСЗІ} --- отримання атестату відповідності КСЗІ України. (2 місяці).
\end{enumerate}

\newpage
\section{Відомості про робочий процес та умови експлуатації}

\subsection{Опис бізнес-процесів (BPMN / FSM)}

\subsubsection{Процес автентифікації клієнта за сертифікатом (chat\_auth.erl)}
Керує процесом перевірки клієнтського сертифіката X.509 при встановленні зв'язку:
\begin{enumerate}
    \item \textbf{ClientConnect}: Клієнт ініціює TCP/TLS з'єднання.
    \item \textbf{RequestCertificate}: Сервер вимагає клієнтський сертифікат (Mutual TLS).
    \item \textbf{VerifyCertificateChain}: Перевірка підпису та ланцюжка сертифікатів через локальний CA.
    \item \textbf{CheckOCSPStatus}: Запит до локального OCSP-сервера для перевірки наявності сертифіката в списку відкликаних.
    \item \textbf{LoadLDAPAttributes}: Читання профілю користувача з LDAP для перевірки ABAC-атрибутів доступу.
    \item \textbf{GrantAccess}: Створення сесії зв'язку, генерація тимчасового токена (AccessToken).
\end{enumerate}

\subsubsection{Процес обміну повідомленнями за протоколом Buddha (chat\_message.erl)}
Керує життєвим циклом передачі повідомлень (CHAT v2):
\begin{enumerate}
    \item \textbf{ReceiveEnvelope}: Прийом DER-кодованого CMS-пакета від відправника.
    \item \textbf{RouteByProfile}: Читання незашифрованого заголовка отримувача та пошук його сесії.
    \item \textbf{StoreTransientQueue}: Збереження зашифрованого конверта в оперативній таблиці Mnesia.
    \item \textbf{DeliverToClient}: Відправка пакета отримувачу при встановленні з'єднання.
    \item \textbf{WaitDeliveryReceipt}: Очікування підтвердження отримання (Ack) від клієнта.
    \item \textbf{DestroyMessage}: Повне видалення повідомлення з оперативної пам'яті сервера.
\end{enumerate}

\subsection{Опис ролей та прав доступу (ABAC)}
Автентифікація користувачів базується на атрибутах LDAP. Приклад Erlang-функції авторизації доступу до каналів зв'язку:
\begin{verbatim}
allow(User, Action, Channel) ->
    UserSecurityLevel = maps:get(clearance_level, User),
    ChannelSecurityLevel = maps:get(required_level, Channel),
    UserDept = maps:get(department, User),
    ChannelDept = maps:get(department, Channel),
    
    HasClearance = (UserSecurityLevel >= ChannelSecurityLevel),
    MatchesDept = (UserDept == ChannelDept),
    
    case {HasClearance, Action} of
        {true, write_message} -> MatchesDept;
        {true, read_message} -> true;
        _ -> false
    end.
\end{verbatim}

\subsection{Специфікація апаратного забезпечення та мережеві порти}
Підсистема «Комунікатор» використовує такі мережеві порти на периферійному сервері:
\begin{center}
\small
\begin{longtable}{p{2cm} p{3.5cm} p{2.5cm} p{5.5cm}}
\toprule
\textbf{Порт} & \textbf{Протокол} & \textbf{Сервіс} & \textbf{Призначення} \\
\midrule
8000 & TCP (WebSocket) & CHAT v2 & Асинхронний обмін DER-повідомленнями \\
4221 & UDP (Multicast) & CHAT v1 / Discovery & Локальне виявлення пристроїв у LAN \\
8443 & TCP (HTTPS) & IAS / EST & Реєстрація клієнтів та перевипуск сертифікатів \\
8089 & TCP (HTTP) & CA / OCSP & Онлайн-перевірка статусу сертифікатів \\
8088 & TCP (HTTP) & CA / CMP & Управління життєвим циклом сертифікатів \\
636 & TCP (LDAPS) & LDAP & Служба каталогів користувачів \\
51820 & UDP & VPN / Тунелі & Захищений WireGuard-канал \\
465 & TCP (SMTPS) & MAIL v3 & Відправка документів (SMTP over TLS) \\
993 & TCP (IMAPS) & MAIL v3 & Отримання документів (IMAP over TLS) \\
\bottomrule
\end{longtable}
\end{center}

\subsection{Архітектурна діаграма взаємодії та зв'язків додатків}

Взаємодія між сервером та клієнтом, а також внутрішня інтеграція Erlang/OTP додатків у межах вузла «Комунікатор» представлена на схемі:

\begin{figure}[ht]
\centering
\begin{tikzpicture}[
    node distance=1.5cm and 2.5cm,
    block/.style={draw=primary!80, fill=primary!5, rectangle, minimum height=2.8em, minimum width=6.2em, rounded corners=4pt, align=center, font=\small, line width=0.8pt},
    db/.style={draw=green!60!black!80, fill=green!5, rectangle, minimum height=2.8em, minimum width=8.0em, rounded corners=4pt, align=center, font=\small, line width=0.8pt},
    client/.style={draw=orange!80, fill=orange!5, rectangle, minimum height=4.5em, minimum width=9.5em, rounded corners=6pt, align=center, font=\small, thick, line width=1.2pt},
    arrow/.style={-{Stealth[scale=1.0]}, thick, draw=primary!80}
]

    % Node definitions
    % Top Row (Backend servers)
    \node[block] (ca)   at (0.5, 0)   {\textbf{CA} \\ (Сертифікати)};
    \node[block] (vpn)  at (3.0, 0)   {\textbf{VPN} \\ (Тунелі)};
    \node[block] (ias)  at (5.5, 0)   {\textbf{IAS} \\ (Авторизація)};
    \node[block] (ldap) at (8.0, 0)   {\textbf{LDAP} \\ (Директорія)};
    \node[block] (chat) at (10.5, 0)  {\textbf{CHAT v2} \\ (Месенджер)};
    \node[block] (mail) at (13.0, 0)  {\textbf{MAIL v3} \\ (Документи)};

    % Bottom Row (Client and Mnesia Database)
    \node[client] (cli) at (3.0, -3.2) {\textbf{SwiftUI Клієнт} \\ (Комунікатор)};
    \node[db] (kvs)     at (10.5, -3.2) {\textbf{Mnesia / KVS} \\ (Сховище)};

    % Boundary box for the Erlang/OTP Node services (excluding the Client)
    \draw[dashed, draw=gray!40, line width=1pt, rounded corners=6pt] (-1.2, -0.8) rectangle (14.7, 1.2);
    \draw[dashed, draw=gray!40, line width=1pt, rounded corners=6pt] (8.5, -4.0) rectangle (12.5, -0.8);
    \node[anchor=north east, font=\footnotesize\bfseries\color{gray}] at (14.6, 1.1) {Вузол Erlang/OTP (Комунікатор)};

    % Connections: Client to Backend servers (clean rays from cli.north)
    \draw[arrow] (cli.north) to[out=135, in=-90] node[left, font=\tiny, pos=0.8] {OCSP} (ca.south);
    \draw[arrow] (cli.north) -- node[right, font=\tiny, pos=0.5] {Tunnel} (vpn.south);
    \draw[arrow] (cli.north) to[out=70, in=-90] node[right, font=\tiny, pos=0.7] {Port 8443} (ias.south);
    \draw[arrow] (cli.north) to[out=50, in=-90] node[right, font=\tiny, pos=0.6] {LDAP} (ldap.south);
    \draw[arrow] (cli.north) to[out=30, in=-90] node[right, font=\tiny, pos=0.5] {Port 8000} (chat.south);
    \draw[arrow] (cli.north) to[out=15, in=-90] node[right, font=\tiny, pos=0.4] {Port 993} (mail.south);

    % Connections: Internal connections to IAS (Auth / Verification)
    \draw[arrow] (ca.north) to[out=40, in=140] node[above, font=\tiny] {EST/CMP} (ias.north);
    \draw[arrow] (vpn.east)   -- (ias.west);
    \draw[arrow] (ldap.west)  -- (ias.east);

    % Connections: Services to Mnesia Database
    \draw[arrow] (chat.south) -- (kvs.north);
    \draw[arrow] (mail.south) to[out=-90, in=35] (kvs.east);

\end{tikzpicture}
\caption{Архітектурна діаграма взаємодії OTP-додатків та мережевих портів}
\label{fig:arch_diag}
\end{figure}

\subsection{Детальна архітектура Комунікатора та інтеграція репозиторіїв}
Для побудови підсистеми використовуються репозиторії екосистеми ERP/1:
\begin{enumerate}
    \item \textbf{Репозиторії екосистеми Erlang/OTP} (шляхи пошуку):
    \begin{itemize}
        \item \texttt{synrc/chat/} --- Сервер CHAT v2 та клієнтські N2O/Nitro залежності.
        \item \texttt{synrc/ca/} --- Локальний сертифікаційний центр.
        \item \texttt{synrc/ldap/} --- Реалізація вбудованого LDAP сервера.
        \item \texttt{zencrypted/ias/} --- Сервер авторизації та ідентифікації.
        \item \texttt{zencrypted/vpn/} та \texttt{../../zencrypted/vpn2/} --- Модулі VPN тунелів.
        \item \texttt{erpuno/mail/} --- Сервер документів MAIL v3.
        \item \texttt{zencrypted/protocol/} --- Опис та специфікації протоколу Buddha.
    \end{itemize}
\end{enumerate}

\newpage
\section{Інформаційне та програмне забезпечення}

\subsection{Інформаційне забезпечення (Схеми та Дані)}

\subsubsection{Визначення структур даних (Erlang Records)}
\begin{verbatim}
-record(chat_session, {
    session_id,                 % UUID сесії (PK)
    client_id,                  % Серійний номер сертифіката X.509
    device_id,                  % Однозначний ідентифікатор пристрою
    ip_address,                 % Поточна IP-адреса
    access_token,               % Токен доступу сесії
    status = active,            % active | expired | revoked
    created_at = 0
}).

-record(chat_roster, {
    roster_id,                  % UUID ростера (PK)
    owner_id,                   % UUID власника
    contacts = []               % Список контактів
}).

-record(ca_certificate, {
    serial_number,              % Серійний номер X.509 (PK)
    subject_dn,                 % Відмітне ім'я суб'єкта
    public_key_der,             % Публічний ключ у DER кодуванні
    not_before = 0,             % Timestamp початку дії
    not_after = 0,              % Timestamp завершення дії
    status = valid              % valid | revoked | expired
}).

-record(vpn_tunnel_state, {
    tunnel_id,                  % UUID тунелю (PK)
    client_ip,                  % Виділена IP-адреса клієнта
    public_key,                 % Публічний ключ WireGuard
    bytes_sent = 0,             % Надіслано байт
    bytes_recv = 0,             % Отримано байт
    last_handshake = 0          % Timestamp останнього рукостискання
}).
\end{verbatim}

\subsubsection{Специфікації ASN.1 для обміну даними}

Нижче наведено специфікацію першого шару для UDP-мультікаст та локальних мереж з файлу \texttt{MESSAGE-v1.asn1}:
\begin{verbatim}
MESSAGE-v1 DEFINITIONS ::= BEGIN

    IMPORTS ContentInfo FROM CryptographicMessageSyntax-2009
            Certificate FROM PKIX1Explicit-2009 ;

    OS  ::= ENUMERATED { apple (1), microsoft (2), google (3),
                         ericsson (4), nokia (5), cisco (6) }
    P2P ::= SEQUENCE { from OCTET STRING, to OCTET STRING }
    MUC ::= SEQUENCE { room OCTET STRING }

    CHATString ::= OCTET STRING

    Feature ::= SEQUENCE { id OCTET STRING, key OCTET STRING, 
                           value OCTET STRING, group OCTET STRING }
    AuthorityMethod ::= ENUMERATED { announce (0), register (1),
                                     authenticate (2), dismiss (3),
                                     renew (4) }
    Authority ::= SEQUENCE {
       session OCTET STRING,
       type AuthorityMethod,
       cert OCTET STRING,
       settings SEQUENCE OF feature Feature }

    File ::= SEQUENCE {
       id OCTET STRING,
       mime OCTET STRING,
       payload ANY,
       parentid OCTET STRING }

    Channel ::= SEQUENCE { channel OCTET STRING }
    MessageType ::= ENUMERATED { sys (1), reply (2), forward (3), 
                                 read (4), edited (5) }
    MessageStatus ::= ENUMERATED { async (1), delete (2), clear (3), 
                                   update (4), edit (5) }
    MessageFeed ::= CHOICE { p2p [0] P2P, muc [1] MUC, chan [2] Channel }
    Message ::= SEQUENCE {
       id CHATString,
       feed MessageFeed,
       from CHATString,
       to CHATString,
       files SEQUENCE OF file File,
       type MessageType,
       link CHOICE { empty [1] NULL, message [2] Message },
       seenby CHATString,
       repliedby CHATString,
       mentioned SEQUENCE OF name CHATString,
       status MessageStatus DEFAULT clear }

    Ack ::= SEQUENCE { table OCTET STRING, id OCTET STRING }

END
\end{verbatim}

\subsection{Програмне забезпечення (Реалізація)}

\subsubsection{chat\_message.erl (DSL процесу обміну повідомленнями)}
\begin{verbatim}
-module(chat_message).
-include("bpe.hrl").
-compile(export_all).

def() ->
    #process{
        name = 'Buddha Protocol Message Delivery',
        beginEvent = 'ReceiveEnvelope',
        endEvent = 'DestroyMessage',
        tasks = [
            #beginEvent{name='ReceiveEnvelope'},
            #serviceTask{name='RouteByProfile', module=?MODULE},
            #serviceTask{name='StoreTransientQueue', module=?MODULE},
            #serviceTask{name='DeliverToClient', module=?MODULE},
            #serviceTask{name='WaitDeliveryReceipt', module=?MODULE},
            #endEvent{name='DestroyMessage'}
        ],
        flows = [
            #sequenceFlow{name='1', source='ReceiveEnvelope', target='RouteByProfile'},
            #sequenceFlow{name='2', source='RouteByProfile', target='StoreTransientQueue'},
            #sequenceFlow{name='3', source='StoreTransientQueue', target='DeliverToClient'},
            #sequenceFlow{name='4', source='DeliverToClient', target='WaitDeliveryReceipt'},
            #sequenceFlow{name='5', source='WaitDeliveryReceipt', target='DestroyMessage'}
        ],
        roles = [operator, system]
    }.

action({serviceTask, 'RouteByProfile'}, Proc) ->
    % Визначення одержувача та активної сесії
    {reply, Proc};
action({serviceTask, 'StoreTransientQueue'}, Proc) ->
    % Запис у транзієнтну Mnesia таблицю
    {reply, Proc};
action({serviceTask, 'DeliverToClient'}, Proc) ->
    % Відправка через TCP сокет
    {reply, Proc};
action(_, Proc) -> {reply, Proc}.
\end{verbatim}

\subsubsection{chat\_nif.c (C99 NIF обгортка для декодування ASN.1 DER)}
Код Erlang NIF модуля мовою C99 для швидкого синтаксичного розбору ASN.1 DER пакетів без Rust:
\begin{verbatim}
#include "erl_nif.h"
#include <string.h>

// Функція розбору ASN.1 TLV (Type-Length-Value) блоку
static ERL_NIF_TERM parse_tlv_nif(ErlNifEnv* env, int argc, const ERL_NIF_TERM argv[]) {
    ErlNifBinary bin;
    if (!enif_inspect_binary(env, argv[0], &bin)) {
        return enif_make_badarg(env);
    }

    if (bin.size < 2) {
        return enif_make_tuple2(env, enif_make_atom(env, "error"), 
                                enif_make_string(env, "too_short", ERL_NIF_LATIN1));
    }

    unsigned char tag = bin.data[0];
    unsigned int length = bin.data[1];
    unsigned int header_len = 2;

    // Спрощена підтримка довгих довжин DER
    if (length & 0x80) {
        unsigned int num_bytes = length & 0x7F;
        if (bin.size < 2 + num_bytes) {
            return enif_make_tuple2(env, enif_make_atom(env, "error"), 
                                    enif_make_string(env, "invalid_length", ERL_NIF_LATIN1));
        }
        length = 0;
        for (unsigned int i = 0; i < num_bytes; i++) {
            length = (length << 8) | bin.data[2 + i];
        }
        header_len = 2 + num_bytes;
    }

    if (bin.size < header_len + length) {
        return enif_make_tuple2(env, enif_make_atom(env, "error"), 
                                enif_make_string(env, "payload_mismatch", ERL_NIF_LATIN1));
    }

    ERL_NIF_TERM type_term = enif_make_uint(env, tag);
    
    unsigned char* payload_data;
    ERL_NIF_TERM payload_term;
    payload_data = enif_make_new_binary(env, length, &payload_term);
    memcpy(payload_data, bin.data + header_len, length);

    return enif_make_tuple3(env, enif_make_atom(env, "ok"), type_term, payload_term);
}

static ErlNifFunc nif_funcs[] = {
    {"parse_tlv", 1, parse_tlv_nif}
};

ERL_NIF_INIT(chat_nif, nif_funcs, NULL, NULL, NULL, NULL)
\end{verbatim}

\newpage
\section{Регламент взаємодії X.422 та криптографічні стандарти}

Регламент взаємодії X.422 визначає дворівневий стандарт обміну:

\subsection{X.422.1: Транспортний рівень}
\begin{itemize}
    \item Взаємодія між клієнтами та брокерами зв'язку CHAT v2 виконується через асинхронні WebSocket з'єднання, захищені двостороннім mTLS 1.3.
    \item Поштовий обмін MAIL v3 реалізує доставку підписаних S/MIME конвертів за стандартними захищеними SMTP/IMAP протоколами.
    \item Для випуску, оновлення та перевірки статусів сертифікатів використовуються стандарти EST (RFC 7030) та OCSP (RFC 6960).
\end{itemize}

\subsection{X.422.2: Криптографічний конверт CMS}
Всі текстові та медіа-повідомлення, а також поштові вкладення перед відправленням проходять обов'язкову криптографічну обробку на боці клієнта:
\begin{enumerate}
    \item \textbf{Формування підпису (Signature)}: створюється CMS-структура SignedData. Обчислення хешу виконується за алгоритмом ДСТУ 7564 (Купина-384) або SHA-384. Електронний підпис обчислюється на еліптичних кривих за стандартом ДСТУ 4145-2002 або ECDSA.
    \item \textbf{Шифрування вмісту (EnvelopedData)}: дані шифруються симетричним шифром ДСТУ 7624 (Калина-256) або AES-GCM-256 на сесійному ключі.
    \item \textbf{Обмін сесійними ключами}: симетричний ключ шифрується на публічному ключі отримувача за алгоритмом еліптичних кривих ECDH (SECP-384 або SECP-571).
\end{enumerate}

\newpage
\section{Вимоги до засобу за ISO/IEC/IEEE 29148}

\subsection{Вимоги до надійності та продуктивності}
\begin{itemize}
    \item Середній показник доступності (SLA) периферійного сервера «Комунікатор» не менше 99.9\%.
    \item Швидкість обробки та маршрутизації повідомлень брокером CHAT v2 — не менше 10 000 повідомлень на секунду при затримці доставки не більше 15 мс.
    \item Пропускна здатність поштової підсистеми MAIL v3 — не менше 500 повідомлень на секунду.
\end{itemize}

\subsection{Вимоги до безпеки та захисту інформації}
\begin{itemize}
    \item Незмінність журналів аудиту, автентифікації користувачів та відкликання сертифікатів.
    \item Заборона збереження та розшифрування закритого тексту (plaintext) повідомлень та поштових вкладень на серверних вузлах.
\end{itemize}

\end{document}
